% Guia do professor — Capítulo 4: Vapor d'Água
\section{Capítulo 4 --- Vapor d'Água}

\subsection*{Visão geral e ritmo}

O dicionário quantitativo da umidade: Clausius--Clapeyron, as sete
variáveis de umidade, LCL, bulbo úmido, água precipitável e o balanço de
água. \textbf{Ritmo sugerido: 2 aulas de 2\,h.}

\begin{itemize}
  \item \textbf{Aula 1} — motivação, dedução e integração de
        Clausius--Clapeyron, a curva $e_s(T)$, o zoológico de variáveis
        (com o exemplo de São Paulo). Fechar com Casos~1--2 (15\,min).
  \item \textbf{Aula 2} — LCL (com dedução da regra dos 125\,m/°C), bulbo
        úmido e limite fisiológico, água precipitável, balanço de água.
        Fechar com Casos~3--4 (10\,min).
\end{itemize}

\subsection*{Objetivos de aprendizagem}

\begin{enumerate}
  \item deduzir e integrar Clausius--Clapeyron; explicar a regra dos
        7\,\%/K e suas consequências climáticas;
  \item converter fluentemente entre $e$, $r$, $q$, $\rho_v$, UR, $T_d$,
        $T_w$ e saber qual usar em cada problema;
  \item explicar por que $r$ é conservada e UR não presta para comparar
        massas de ar;
  \item calcular o LCL pela regra e pelo método exato; ler a base do
        cumulus como higrômetro;
  \item calcular água precipitável de uma sondagem e montar o balanço
        euleriano de água.
\end{enumerate}

\subsection*{Pré-requisitos a reativar}

Gás ideal por espécie (pressões parciais, Cap.~1); $\Gamma_d$ e conservação
de $\theta$ (Cap.~3); potencial de Gibbs é citado mas a dedução de
Clausius--Clapeyron pode ser dada ``por decreto'' com a EDO como ponto de
partida — físicos aceitam bem.

\subsection*{Armadilhas conceituais frequentes}

\begin{itemize}
  \item \textbf{``Ar quente segura mais água''}: linguagem de esponja
        esconde a física — é a \emph{pressão de equilíbrio} $e_s(T)$ que
        cresce; o ar não ``segura'' nada. Vale corrigir explicitamente.
  \item \textbf{UR como medida absoluta}: 90\,\% de UR no inverno pode ser
        menos vapor que 40\,\% no verão. O exemplo numérico cura.
  \item \textbf{$T_d$ vs.\ $T_w$}: alunos trocam; fixar
        $T_d \le T_w \le T$ com o caso saturado degenerando tudo.
  \item \textbf{Fator $\varepsilon$}: esquecem o 0,622 e erram $r$ por
        60\,\%; lembrar a origem molecular (18/29).
  \item \textbf{LCL para ar não aquecido pela superfície}: a regra dos
        125\,m/°C usa $T$ e $T_d$ \emph{da parcela que sobe}; para
        estratiforme o levantamento é outro (Cap.~5).
\end{itemize}

\subsection*{Demonstração de baixo custo}

Psicrômetro artesanal: dois termômetros idênticos, gaze molhada num deles,
abanar 1 minuto. Medir $T$ e $T_w$, calcular $r$ e UR em sala com a equação
psicrométrica (é o exercício N2 ao vivo). Alternativa: lata gelada suando —
a película é o orvalho, a temperatura da lata no início da condensação
$\approx T_d$ do ar da sala.

\subsection*{Problemas sugeridos do livro (seção de exercícios do Cap.~4)}

Aquecimento: cálculo de $e_s$ para várias $T$; conversões $e \to r \to$ UR;
$T_d$ e LCL. Aprofundamento: bulbo úmido e psicrômetro; água precipitável
de perfis dados. Síntese: ``por que o deserto tem a maior amplitude térmica
diária?'' (liga UR, $I\!\downarrow$ do Cap.~2 e Bowen do Cap.~3).
\emph{Confira a numeração na sua cópia (v1.02b).}

\subsection*{Uso do notebook \texttt{cap04\_vapor.ipynb}}

\begin{center}
\begin{tabular}{lll}
\hline
Caso & Conteúdo & Momento sugerido \\
\hline
1 & EDO de Clausius--Clapeyron × fórmula × MetPy; gelo × líquido & Aula 1, após CC \\
2 & Carta psicrométrica $r$--$T$ com curvas de UR & Aula 1, após variáveis \\
3 & LCL exato × regra; base de nuvem ao longo do dia & Aula 2, após LCL \\
4 & Água precipitável; escala de 7\,\%/K & fim da Aula 2 \\
\hline
\end{tabular}
\end{center}

O painel gelo × líquido do Caso~1 planta a semente do processo de Bergeron
(Cap.~7); o Caso~4 fecha com a conexão climática (chuvas extremas).

\subsection*{Exercícios complementares e soluções}

Nas notas: T1--T5 (integração de CC com $L_v$ variável, regra dos 7\,\%/K,
dedução de $r(\varepsilon)$, dedução da regra do LCL, água precipitável
analítica) e N1--N4 (células no notebook). Soluções em
\texttt{cap04\_solucoes.ipynb} (\emph{não distribuir}): T1 e T4 com
\texttt{sympy}, N2 com \texttt{brentq} validado contra o MetPy.
Pares teoria--numérica: T1$\leftrightarrow$N1 e T2$\leftrightarrow$N4.

\subsection*{Conexões com o resto do curso}

$e_s$ sobre gelo × líquido $\to$ Bergeron (Cap.~7); LCL $\to$ base de nuvem
e diagramas (Caps.~5--6); $r$ conservada + $\theta$ $\to$ diagramas
termodinâmicos (Cap.~5); água precipitável $\to$ rios atmosféricos e chuvas
extremas (Caps.~13, 21); balanço de água $\to$ ciclo hidrológico (Cap.~21).
